\section{Discussion and Limitations}\label{sec:discussion}

\paragraph{The single-component boundary.} The certificate helps only when the process
model has more than one component. A single process whose cases share a resource pool
is one component, and there our pipeline is exactly PPO. This is common: in the BPI
Challenge 2012 loan-application log~\cite{bpic2012}, a median case is handled by four
different resources, and the resources form one connected group even without the
validation team. The target of our method is therefore multi-process or multi-team
control: several processes with dedicated teams under one indicator.

\paragraph{Per-process learners.} When the processes are independent, separate learners
also avoid the other processes' noise, and they capture most of the gain in credit. What
the certificate adds is a single shared policy that learns faster early,
postponement without freezing, an automatic decomposition that is sound by
construction, and the correct merging of processes that turn out to be coupled.

\paragraph{Granularity of postponement.} A component waits as a whole until one of its
events; the agent can assign some resources first and then wait, but it cannot keep one
resource idle while another keeps receiving work in the same decision round.
\todo{check whether this needs qualification}

\paragraph{Synthetic benchmark.} The insurer is designed, not mined from data. Its
sub-processes reuse validated models, and the waiting regime is grounded in queueing
theory, but a multi-organisation case from real data, such as the five municipalities of
the BPI Challenge 2015, would strengthen the external validity. \todo{BPIC 2015 if the
staff of the municipalities turn out to be disjoint}
